\mxnewpage\section*{EV1-A -- DS commun}\label{EV1-A}
\addcontentsline{toc}{section}{EV1-A -- DS commun}
\textbf{Sujet réservé au professeur}. Chapitres 1 à 4.\par Durée : 2 h. Sans calculatrice. Note sur 20. 14 novembre 2026.
\subsection*{Partie 1 -- Automatismes (6 points)}
Douze questions, \num{0,5} point chacune. Une seule réponse correcte ; aucune justification demandée. Une réponse incorrecte ou absente rapporte 0, sans retrait de points.
\begin{enumerate}
\item Calculer $5+2\times3$.\par A. $14$\quad B. $11$\quad C. $10$\quad D. $21$
\item Calculer $\frac12+\frac13$.\par A. $\frac{5}{6}$\quad B. $\frac{2}{5}$\quad C. $\frac{5}{3}$\quad D. $\frac{1}{6}$
\item Calculer $\dfrac{2^5}{2^{3}}$.\par A. $2$\quad B. $1$\quad C. $10$\quad D. $4$
\item Une valeur de 50 augmente de \pc{20}. Valeur finale ?\par A. $10$\quad B. $60$\quad C. $70$\quad D. $40$
\item Une valeur de 80 est obtenue après une baisse de \pc{20}. Valeur initiale ?\par A. $100$\quad B. $60$\quad C. $64$\quad D. $96$
\item Hausse de \pc{25} puis baisse de \pc{20} : coefficient global ?\par A. $\frac{9}{10}$\quad B. $1$\quad C. $\frac{21}{20}$\quad D. $\frac{4}{5}$
\item La solution de $2x+5=0$ vaut\par A. $10$\quad B. $\frac{-5}{2}$\quad C. $-10$\quad D. $\frac{5}{2}$
\item La pente entre $\pt{0;5}$ et $\pt{2;11}$ vaut\par A. $6$\quad B. $3$\quad C. $2$\quad D. $-3$
\item La moyenne de $5, 7, 9$ vaut\par A. $8$\quad B. $21$\quad C. $5$\quad D. $7$
\item Un secteur de $90^\circ$ représente quelle fraction du total ?\par A. $\frac{1}{3}$\quad B. $\frac{1}{2}$\quad C. $\frac{1}{4}$\quad D. $\frac{3}{4}$
\item Une probabilité vaut $\dfrac{2}{5}$. Son complément vaut\par A. $\frac{3}{5}$\quad B. $\frac{2}{5}$\quad C. $\frac{5}{2}$\quad D. $\frac{1}{5}$
\item Pour $f(x)=x^2+5$, $f(-2)$ vaut\par A. $9$\quad B. $7$\quad C. $1$\quad D. $20$
\end{enumerate}
\mxnewpage\subsection*{Partie 2 -- Exercices indépendants (14 points)}
Toute démarche doit être justifiée. Les résultats exacts sont attendus sauf indication contraire.
\mxneed{12}
\subsection*{Exercice 1 -- Réserve décroissante (5 points)}
Une réserve contient initialement 192 L d’eau. À chaque étape, on retire la moitié de l’eau encore présente, sans ajouter d’eau. On note $u_n$ le volume, en litres, restant après $n$ étapes ; le rang 0 correspond à l’état initial.
\begin{enumerate}
\item Définir la suite avec son premier terme et sa raison. \hfill(1 pt)
\item Calculer les volumes des rangs 0 à 4 et leur somme. \hfill(2 pts)
\item Déterminer le premier rang où le volume est strictement inférieur à 10 L ; décrire une boucle de seuil. \hfill(2 pts)
\end{enumerate}
\mxneed{12}
\subsection*{Exercice 2 -- Paramètre et tangente (5 points)}
Pour un réel $m$ fixé, on définit sur $\R$ la fonction $f_m(x)=x^2-2mx+4$. Le nombre $m$ est un paramètre. Les deux premières questions portent sur toutes ses valeurs ; la dernière concerne uniquement $m=1$.
\begin{enumerate}
\item Calculer le discriminant du polynôme $f_m$ en fonction de $m$. \hfill(1 pt)
\item Déterminer, selon $m$, le nombre de solutions réelles de $f_m(x)=0$. \hfill(2 pts)
\item On fixe maintenant $m=1$. Calculer le taux de variation entre 2 et $2+h$, pour $h\ne0$, puis l’équation de la tangente au point d’abscisse 2. \hfill(2 pts)
\end{enumerate}
\mxneed{12}
\subsection*{Exercice 3 -- Locus par le milieu (4 points)}
Deux points fixes $A$ et $B$ du plan sont distants de 8 unités. On note $I$ le milieu du segment $[AB]$ et $M$ un point variable du plan. On cherche à décrire des ensembles de points à partir du produit scalaire $\vect{MA}\cdot\vect{MB}$.
\begin{enumerate}
\item Exprimer $\vect{MA}\cdot\vect{MB}$ en fonction de $MI$, en justifiant l’identité utilisée. \hfill(1 pt)
\item Décrire l’ensemble des points $M$ pour lesquels ce produit vaut 9. \hfill(1 pt)
\item Décrire successivement les ensembles obtenus lorsque le produit vaut $-16$, puis $-20$. \hfill(2 pts)
\end{enumerate}
